\chapter{Score matching}
\label{ch:score-matching}

The paper's networks are trained to approximate the score of the noised law, and its
structured estimator is compared with them through the score. This chapter explains
what ``learning a score'' means, derives the two objectives that make it possible
without ever knowing the score (Hyv\"arinen's implicit and Vincent's denoising score
matching), works out what each becomes under the Ornstein--Uhlenbeck channel of
\cref{ch:ou-channel-tweedie}, and then follows the code: the two output
parameterisations, the three ways the divergence term is computed, the minibatch
estimate of the loss, the estimand the paper reports, and the three information
interfaces at which the comparison is run.

\section{What learning a score means}
\label{sec:sm:meaning}

A \term{score model} is a family of vector fields $s_\params(\vx,t)\colon\Reals^{\nsites}\times(0,\infty)\to\Reals^{\nsites}$
indexed by parameters $\params$ (the weights of a network). The target is the score
$\score(\vx,t)=\nabla_{\vx}\log\Pt(\vx)$ of the noised law, and the natural measure of
success is the \term{explicit score-matching} loss at one level,
\begin{equation}
  J_{\mathrm{ESM}}(s)=\E_{\Pt}\bigl|s(\vx_t)-\nabla\log\Pt(\vx_t)\bigr|^2 .
  \label{eq:sm:esm}
\end{equation}
The difficulty is that $\nabla\log\Pt$ is unknown: for real data one has samples of
$\Pt$ (or of $\Pz$, which one can noise) and nothing else. In the paper's laboratory
$\score$ is known exactly (Section 2), which is what allows $J_{\mathrm{ESM}}$ to be
measured; but the networks are still trained as they would be in practice, by one of
the two surrogates below, which have the same minimiser as \cref{eq:sm:esm} and can be
estimated from samples. In both, $t$ is fixed for the moment; the paper's schedule
enters in \cref{sec:sm:batch}.

\section{Hyv\"arinen's implicit score matching}
\label{sec:sm:ism}

\begin{theorem}[Implicit score matching]
\label{thm:sm:ism}
Let $p$ be a continuously differentiable density on $\Reals^{\nsites}$ and let
$s\colon\Reals^{\nsites}\to\Reals^{\nsites}$ be continuously differentiable, with
$\E_p|s|^2<\infty$, $\E_p|\nabla\log p|^2<\infty$, $\E_p|\operatorname{div}s|<\infty$,
and $p(\vx)\,|s(\vx)|\to0$ as $|\vx|\to\infty$. Then
\begin{equation}
  \E_p\bigl|s-\nabla\log p\bigr|^2
  =2\,\E_p\Bigl[\tfrac12|s|^2+\operatorname{div}s\Bigr]+\E_p|\nabla\log p|^2 .
  \label{eq:sm:ism}
\end{equation}
The last term does not depend on $s$, so the \term{implicit score-matching} objective
$J_{\mathrm{ISM}}(s)=\E_p[\tfrac12|s|^2+\operatorname{div}s]$ has the same minimiser as
$J_{\mathrm{ESM}}$, and $J_{\mathrm{ISM}}$ involves only $s$ and samples of $p$.
\end{theorem}

\begin{proof}
Scalar case first. Expand the square:
\[
  \int p\,(s-(\log p)')^2\dd x=\int p\,s^2\dd x-2\int p\,s\,\frac{p'}{p}\dd x+\int p\,((\log p)')^2\dd x .
\]
In the middle term $p\cdot p'/p=p'$, so it is $-2\int s\,p'\dd x$, and integration by
parts on $[-R,R]$ gives
\[
  \int_{-R}^{R}s\,p'\dd x=\bigl[s\,p\bigr]_{-R}^{R}-\int_{-R}^{R}s'\,p\,\dd x .
\]
The boundary term tends to zero as $R\to\infty$ by the decay hypothesis, and the
remaining integral converges by the integrability of $s'$ against $p$. Hence
$\int p\,(s-(\log p)')^2=\int p\,s^2+2\int p\,s'+\int p\,((\log p)')^2$, which is
\cref{eq:sm:ism} with $\operatorname{div}s=s'$. In $\nsites$ dimensions the middle term is
$-2\int\langle s,\nabla p\rangle\dd\vx$, and the divergence theorem on the ball $B_R$
gives $\int_{B_R}\langle s,\nabla p\rangle=\oint_{\partial B_R}p\,\langle s,\vect{n}\rangle\dd\sigma-\int_{B_R}p\operatorname{div}s$;
the surface term is bounded by the area of the sphere times $\sup_{|\vx|=R}p|s|$, which
tends to zero by hypothesis (the hypothesis is to be read as decay fast enough for
this product to vanish).
\end{proof}

Two remarks. First, at the true score $J_{\mathrm{ISM}}(\score)=-\tfrac12\E|\score|^2$,
i.e.\ $\E_p[\operatorname{div}\nabla\log p]=-\E_p|\nabla\log p|^2$; combined with
\cref{prop:ou:tweedie2} this reproduces \cref{prop:ou:energy}. Second, the price of
avoiding $\nabla\log p$ is the divergence, a sum of $\nsites$ diagonal entries of the
Jacobian of $s$, which is what \cref{sec:sm:divergence} is about.

\section{Vincent's denoising score matching}
\label{sec:sm:dsm}

Suppose the samples of $p$ arise by corrupting clean samples: $\va\sim\Pz$ and then
$\vx\sim p(\cdot\mid\va)$, so that $p(\vx)=\int\Pz(\dd\va)\,p(\vx\mid\va)$. The
conditional score $\nabla_{\vx}\log p(\vx\mid\va)$ is known in closed form whenever the
corruption is.

\begin{theorem}[Denoising score matching]
\label{thm:sm:dsm}
Let $p(\vx\mid\va)$ be differentiable in $\vx$ with $\E|\nabla_{\vx}\log p(\vx\mid\va)|^2<\infty$
under the joint law, and let $\E_p|s|^2<\infty$. Then
\begin{equation}
  J_{\mathrm{DSM}}(s):=\E_{(\va,\vx)}\bigl|s(\vx)-\nabla_{\vx}\log p(\vx\mid\va)\bigr|^2
  =J_{\mathrm{ESM}}(s)+\E_{(\va,\vx)}\bigl|\nabla_{\vx}\log p(\vx\mid\va)-\nabla_{\vx}\log p(\vx)\bigr|^2 .
  \label{eq:sm:dsm}
\end{equation}
The second term does not depend on $s$, so $J_{\mathrm{DSM}}$ and $J_{\mathrm{ESM}}$
have the same minimiser, and $J_{\mathrm{DSM}}$ needs only pairs $(\va,\vx)$.
\end{theorem}

\begin{proof}
The key identity is that the marginal score is the conditional expectation of the
conditional score,
\begin{equation}
  \E\bigl[\nabla_{\vx}\log p(\vx\mid\va)\;\big|\;\vx\bigr]=\nabla_{\vx}\log p(\vx),
  \label{eq:sm:tower}
\end{equation}
because, with the posterior $p(\va\mid\vx)=\Pz(\va)p(\vx\mid\va)/p(\vx)$,
\[
  \int p(\va\mid\vx)\,\nabla_{\vx}\log p(\vx\mid\va)\dd\va
  =\frac{1}{p(\vx)}\int\Pz(\va)\,\nabla_{\vx}p(\vx\mid\va)\dd\va
  =\frac{\nabla_{\vx}p(\vx)}{p(\vx)} .
\]
Now write $Y=\nabla_{\vx}\log p(\vx\mid\va)$ and $S=\nabla\log p(\vx)=\E[Y\mid\vx]$. For
any function $s$ of $\vx$, the orthogonality argument of \cref{prop:ou:mmse} applied
to $Y$ gives $\E|s(\vx)-Y|^2=\E|s(\vx)-S|^2+\E|S-Y|^2$, which is \cref{eq:sm:dsm}.
\end{proof}

The regression reading of the proof is worth stating on its own: $J_{\mathrm{DSM}}$ is
a least-squares regression of the random target $Y$ on the input $\vx$, and the
least-squares regression function is the conditional expectation, which by
\cref{eq:sm:tower} is the score. The irreducible part of the loss is the conditional
variance of the target.

\begin{corollary}[The channel]
\label{cor:sm:channel}
Under $\vx_t=e^{-t}\va+\sqrt{\Dt}\vz$, the conditional law is
$\normal{e^{-t}\va}{\Dt\Id}$ and the target is
\begin{equation}
  \nabla_{\vx}\log p(\vx_t\mid\va)=-\frac{\vx_t-e^{-t}\va}{\Dt}=-\frac{\vz}{\sqrt{\Dt}} .
  \label{eq:sm:target}
\end{equation}
Identity \cref{eq:sm:tower} is Tweedie's identity \cref{eq:ou:tweedie} in another
guise, and the irreducible loss is
$\E|{-\vz/\sqrt{\Dt}}-\score|^2=(e^{-2t}/\Dt^2)\,\mathrm{MMSE}_t$ by
\cref{prop:ou:energy}.
\end{corollary}

The target's variance is $\nsites/\Dt$ per example, which diverges as $t\to0$; a network
asked to output the score directly would have to produce values of size $1/\sqrt{\Dt}$
at the lowest levels. The parameterisations of the next section are the standard way to
avoid that.

\section{The two output parameterisations}
\label{sec:sm:param}

Every network of the paper maps a noised sequence and a level to an output
$f_\params(\vx,t)\in\Reals^{\nsites}$, and the output is converted to a posterior mean
and then to a score by \cref{thm:ou:tweedie}. The code supports two conventions;
they are implemented by \code{Denoiser.posterior\_mean}, \code{\_score\_from\_out}
and \code{dsm\_loss} in \code{src/\allowbreak markov\_score/\allowbreak nn\_torch.py}.

\begin{definition}[\code{eps} and \code{x0}]
\label{def:sm:param}
With $\alpha=e^{-t}$:
\begin{itemize}
\item \code{eps}: the output predicts the noise, target $\vz$; the posterior mean is
  $\widehat m=(\vx-\sqrt{\Dt}\,f)/\alpha$ and the score is
  $\widehat s=-(\vx-\alpha\widehat m)/\Dt=-f/\sqrt{\Dt}$.
\item \code{x0}: the output predicts the clean sequence, target $\va$; the posterior
  mean is $\widehat m=f$ and the score is $\widehat s=-(\vx-\alpha f)/\Dt$.
\end{itemize}
The training loss in both cases is the mean squared distance between output and
target over the batch, $\E|f-\vz|^2$ or $\E|f-\va|^2$.
\end{definition}

\begin{proposition}[Both losses are weighted denoising score matching]
\label{prop:sm:weights}
With $s^\ast=-(\vx-\alpha\va)/\Dt$ the target of \cref{eq:sm:target},
\[
  |f-\vz|^2=\Dt\,|\widehat s-s^\ast|^2\ \ (\code{eps}),\qquad
  |f-\va|^2=\frac{\Dt^2}{\alpha^2}\,|\widehat s-s^\ast|^2\ \ (\code{x0}).
\]
Hence the \code{eps} loss is $J_{\mathrm{DSM}}$ weighted by $\Dt$, and the \code{x0} loss
is $J_{\mathrm{DSM}}$ weighted by $\Dt^2/\alpha^2=\Dt/\mathrm{SNR}_t$. Both have the
score as their minimiser at every fixed $t$.
\end{proposition}

\begin{proof}
From $\widehat s=-(\vx-\alpha\widehat m)/\Dt$ and $s^\ast=-(\vx-\alpha\va)/\Dt$,
$\widehat s-s^\ast=\alpha(\widehat m-\va)/\Dt$, so $|\widehat m-\va|^2=(\Dt^2/\alpha^2)|\widehat s-s^\ast|^2$;
this is the \code{x0} case since $\widehat m=f$. In the \code{eps} case
$\widehat m-\va=(\vx-\sqrt{\Dt}f)/\alpha-(\vx-\sqrt{\Dt}\vz)/\alpha=-\sqrt{\Dt}(f-\vz)/\alpha$,
so $|f-\vz|^2=(\alpha^2/\Dt)|\widehat m-\va|^2=\Dt|\widehat s-s^\ast|^2$.
\end{proof}

The weights matter once levels are mixed (\cref{sec:sm:batch}): \code{eps} weights the
levels almost equally in score units at high noise ($\Dt\approx1$) and suppresses the
lowest levels ($\Dt\approx2t$); \code{x0} suppresses them more strongly. The paper
trains both and lets the screen choose (Appendix C, ``both output
parameterisations''). The conversion $\widehat s=-f/\sqrt{\Dt}$ for \code{eps} is the
reason a network output of order one corresponds to a score of order $1/\sqrt{\Dt}$,
as the target requires.

\section{Computing the divergence}
\label{sec:sm:divergence}

Implicit score matching needs $\operatorname{div}\widehat s=\sum_i\partial\widehat s_i/\partial x_i=\tr J$
for every example in a batch, where $J=\partial\widehat s/\partial\vx$ is the
$\nsites\times\nsites$ Jacobian of the network's score. Automatic differentiation
gives Jacobian-transpose-vector products cheaply: for any $\vect{v}$, one backward pass
through $\langle\widehat s,\vect{v}\rangle$ returns $J^{\!\top}\vect{v}$, since
$\partial_j\sum_i\widehat s_i v_i=\sum_i v_iJ_{ij}$. The function \code{divergence} in
\code{src/markov\_score/nn\_torch.py} offers three ways to turn such products into a
trace.

\subsection{The exact trace by $\nsites$ backward passes}
Taking $\vect{v}=\vect{e}_i$ returns the $i$-th row of $J$, whose $i$-th entry is
$J_{ii}$. Summing over $i$ costs $\nsites$ backward passes per batch (the
\code{"exact"} branch: for each $i$, the gradient of $\sum_{\text{batch}}\widehat s_i$
with respect to $\vx$, read at coordinate $i$). At $\nsites=32$ this is thirty-two
passes for every training step, and it is what the frozen matched arm used for the
window head (Appendix C: ``computed exactly for the window head'').

\subsection{Hutchinson's estimator}
\begin{proposition}[Hutchinson]
\label{prop:sm:hutchinson}
Let $\vect{v}$ be a random vector with $\E[\vect{v}\vect{v}^{\!\top}]=\Id$, independent
of $J$. Then $\E[\vect{v}^{\!\top}J\vect{v}]=\tr J$. If $\vect{v}\sim\normal{0}{\Id}$,
\[
  \Var(\vect{v}^{\!\top}J\vect{v})=\Frob{J}^2+\tr(J^2)=\Frob{J}^2+\langle J,J^{\!\top}\rangle,
\]
which equals $2\Frob{J}^2$ when $J$ is symmetric.
\end{proposition}

\begin{proof}
$\E[\vect{v}^{\!\top}J\vect{v}]=\sum_{ij}J_{ij}\E[v_iv_j]=\sum_iJ_{ii}$. For the variance
with Gaussian $\vect{v}$, use $\E[v_iv_jv_kv_l]=\delta_{ij}\delta_{kl}+\delta_{ik}\delta_{jl}+\delta_{il}\delta_{jk}$:
\[
  \E\bigl[(\vect{v}^{\!\top}J\vect{v})^2\bigr]=\sum_{ijkl}J_{ij}J_{kl}\E[v_iv_jv_kv_l]
  =(\tr J)^2+\sum_{ij}J_{ij}^2+\sum_{ij}J_{ij}J_{ji},
\]
and subtract $(\tr J)^2$.
\end{proof}

The score of a density has a symmetric Jacobian (\cref{prop:ou:tweedie2}); a
network's need not, but is close to symmetric when it is close to a score. With one
standard-normal probe per example per step (the \code{"hutchinson"} branch, used by the
frozen matched arm for the transformer), the loss estimate is unbiased but its variance
adds to the minibatch noise of \cref{sec:sm:batch}. Hutchinson (1990) used Rademacher
probes ($v_i=\pm1$), for which the diagonal contributes no variance and
$\Var=2\sum_{i\ne j}J_{ij}J_{ji}$, at most the Gaussian value; the code uses Gaussian
probes. For scale: if $J=\Id_{32}$ a Gaussian probe estimates the trace $32$ with
standard deviation $\sqrt{2\cdot32}=8$.

\subsection{Coloured probes for a banded Jacobian}
\label{sec:sm:coloured}

The window head of the paper computes its output at site $i$ from the inputs at sites
$|j-i|\le r$ only (radius $r$; \code{receptive\_radius} in the code). Its Jacobian is
therefore \term{banded}: $J_{ij}=0$ for $|i-j|>r$. For such a map the exact trace
costs $2r+1$ backward passes rather than $\nsites$.

\begin{proposition}[Coloured probes]
\label{prop:sm:coloured}
Let $J_{ij}=0$ whenever $|i-j|>r$, put $p=2r+1$, and for $k=0,\dots,p-1$ let
$\vect{v}_k\in\{0,1\}^{\nsites}$ have entries $(\vect{v}_k)_j=\mathbf 1\{j\equiv k\pmod p\}$.
Then $\sum_{k=0}^{p-1}\vect{v}_k^{\!\top}J\vect{v}_k=\tr J$ exactly.
\end{proposition}

\begin{proof}
$\vect{v}_k^{\!\top}J\vect{v}_k=\sum_{i\equiv k}\sum_{j\equiv k}J_{ij}$. If $i\equiv j\pmod p$
and $i\ne j$ then $|i-j|\ge p=2r+1>r$, so $J_{ij}=0$. Only $i=j$ survives:
$\vect{v}_k^{\!\top}J\vect{v}_k=\sum_{i\equiv k}J_{ii}$. Every $i$ has exactly one
residue class, so the sum over $k$ is $\sum_iJ_{ii}$.
\end{proof}

This is the colouring of a banded Jacobian of Curtis, Powell and Reid (1974): two
columns may share a probe when no row has a nonzero in both, and the residue classes
modulo $2r+1$ are such a partition. The code (\code{"coloured"} branch) computes, for
each $k$, one backward pass $J^{\!\top}\vect{v}_k$ and then the inner product with
$\vect{v}_k$, exactly as for Hutchinson but with $p$ deterministic probes and no
variance. Zero padding at the ends of the chain only removes dependencies, so the
boundary sites are covered. The branch refuses a model that declares no receptive
radius, since for a dense Jacobian the cross terms would not vanish. The identity is
checked against the $\nsites$-pass trace by the test
\code{test\_coloured\_trace\_\allowbreak equals\_exact\_trace\_\allowbreak for\_the\_window\_head}
in \code{tests/test\_nn\_torch.py}. For the radius-2 head this is five passes in place
of thirty-two.

\begin{example}[Seven sites, radius one]
\label{ex:sm:coloured}
With $\nsites=7$, $r=1$ and $p=3$ the probes are $\vect{v}_0=\vect{e}_1+\vect{e}_4+\vect{e}_7$,
$\vect{v}_1=\vect{e}_2+\vect{e}_5$ and $\vect{v}_2=\vect{e}_3+\vect{e}_6$. For a
tridiagonal $J$,
$\vect{v}_0^{\!\top}J\vect{v}_0=J_{11}+J_{44}+J_{77}+(J_{14}+J_{41}+J_{47}+J_{74}+J_{17}+J_{71})$
and every bracketed entry is zero because its indices differ by at least three. Adding
the three probes gives $J_{11}+\dots+J_{77}$.
\end{example}

\section{Why the loss is estimated per batch}
\label{sec:sm:batch}

The population objective is an expectation; training minimises it by stochastic
gradient descent on unbiased minibatch estimates. Under the paper's recipe
(\code{TrainConfig} and \code{train} in \code{src/markov\_score/nn\_torch.py}) one step
draws a batch of $64$ examples, and for each example a level $t$ uniformly from the
twelve of \cref{tab:ou:levels} and a fresh corruption $\vz$; the loss is the mean of the
per-example losses; the gradient is taken; its norm is clipped at $1$; and Adam takes a
step at learning rate $10^{-3}$ or $3\times10^{-3}$. Because $t$ and $\vz$ are drawn
independently of $\params$, the batch mean is an unbiased estimate of
\begin{equation}
  \bar J(\params)=\frac1{12}\sum_{k=1}^{12}w(t_k)\,J^{(t_k)}_{\mathrm{DSM}}(s_\params(\cdot,t_k))
  \label{eq:sm:pooled-objective}
\end{equation}
with the weights $w$ of \cref{prop:sm:weights}, and its gradient is an unbiased estimate
of $\nabla\bar J$. The minimiser of \cref{eq:sm:pooled-objective} over all vector
fields is the score at every level simultaneously, since the levels do not compete
(the field may depend on $t$); with a finite network they do compete, and the weights
decide the compromise. Checkpoints are saved on a ladder up to $32{,}000$ steps and one
is selected afterwards on a validation bundle; the selection rule is part of the
interface (\cref{sec:sm:interfaces}). For implicit score matching the same loop draws
$64$ examples from the fixed set of noised views, computes $\tfrac12|\widehat s|^2+\operatorname{div}\widehat s$
per example with one of the three divergence branches, and averages; with a Hutchinson
probe the estimate is still unbiased, since the probe is drawn independently of
$\params$, at the cost of extra variance.

\section{The estimand: schedule-pooled relative score error}
\label{sec:sm:estimand}

\begin{definition}[The estimand]
\label{def:sm:estimand}
The paper's headline number for any estimator $\widehat s$ is the
\term{schedule-pooled relative score error} (\code{experiments/protocol.py},
\code{estimand}; Appendix C, ``Estimand''),
\begin{equation}
  \mathrm{err}(\widehat s)=\Biggl(\frac{\sum_{t}\sum_{j}\sum_{i}\bigl(\widehat s_t(\vx_j)_i-s_t(\vx_j)_i\bigr)^2}
                                   {\sum_{t}\sum_{j}\sum_{i}s_t(\vx_j)_i^2}\Biggr)^{1/2},
  \label{eq:sm:estimand}
\end{equation}
where $t$ runs over the twelve levels, $j$ over the sequences of a test bundle noised
at that level, $i$ over the sites of a region, and $s_t$ is the tree-exact reference
score. Both sums are formed before the square root.
\end{definition}

Three features of \cref{eq:sm:estimand} matter for reading the paper. It is
\emph{relative}: an estimator that outputs zero scores $1$, and the isotropic score
$-\vx$ and the LMMSE twin give the two reference values of the paper's oracle ladder.
It is \emph{pooled} rather than averaged over levels: the numerator and the denominator
are energies summed over the schedule, so a level contributes in proportion to its
score energy, which by \cref{prop:ou:energy} is close to $\nsites/\Dt$ per sequence at
low noise. Taking $1/\Dt$ as a proxy for the per-site energy, the twelve levels of
\cref{tab:ou:levels} carry the weights $10.5, 7.4, 5.3, 3.8, 2.8, 2.1, 1.65, 1.35,
1.16, 1.06, 1.02, 1.00$, whose sum is $39.1$; the lowest four levels carry about
$69\%$ of the pooled denominator and the highest four about $11\%$. A pooled error is
therefore a low-noise-dominated summary, which the paper says in words (Section 3,
``What this says about score learning'') and which explains why the fourth-order
Edgeworth correction, exact at high noise, is worse than the twin when pooled
(\cref{ch:attenuation-theorems}). Finally it is computed on \emph{regions}: all
$32$ sites (the headline), the bulk, sites $5$ to $28$ (four trimmed at each end, used
for every selection so that boundary effects do not drive the choice of checkpoint),
and the centre site (a diagnostic). Ratios between arms are formed within seed and
averaged, and intervals are percentile bootstraps over seeds (Appendix C).

\section{The three information interfaces}
\label{sec:sm:interfaces}

The paper's Section 5 declares three interfaces and reports the structured
estimator against the networks at each. The definitions below follow
\code{experiments/e03\_baselines.py} line by line.

\begin{definition}[Interface A: the standard denoising recipe]
\label{def:sm:A}
The network sees the $\nseq$ clean training sequences. At every step it draws a batch,
a level per example from the twelve, and a fresh corruption per example, and minimises
the \code{eps} or \code{x0} denoising loss (\code{\_train\_select} with
\code{loss="dsm"}; \code{\_draw\_batch} in \code{nn\_torch.py}). The checkpoint is
selected on a validation bundle of $1024$ sequences, scored by
\cref{eq:sm:estimand} on the bulk sites against the tree-exact reference under the
true kernel.
\end{definition}

\begin{definition}[Interface B: clean targets, one fixed corruption]
\label{def:sm:B}
Each training sequence is corrupted once, at one level, with the sequences split
evenly across the twelve levels, by the same function \code{\_one\_\allowbreak view\_\allowbreak groups}
that builds the structured estimator's data. The pair $(\va,\vx)$ is kept and never
redrawn, and the target $\vz=(\vx-e^{-t}\va)/\sqrt{\Dt}$ (or $\va$) is regressed by the
same denoising loss on batches drawn from the fixed pairs; the training loop is
\code{\_train\_\allowbreak select\_\allowbreak fixed\_\allowbreak corruption}. The
checkpoint is selected as in interface A, on the validation bundle against the true
score.
\end{definition}

\begin{definition}[Interface C: one noised view per sequence, no oracle]
\label{def:sm:C}
The network sees only the noised views $(\vx^{\mu},t_\mu)$ of interface B, without the
clean sequences: the same one-view bundle the structured estimator sees. It is trained
by implicit score matching (\code{loss="ism"}), with the divergence computed exactly
for the window head and by Hutchinson's estimator for the transformer. In E03 the
exact trace of the window head is the $\nsites$-pass branch
(\code{hutchinson=(arch != "window")} with no \code{trace} override); in E14, whose
network arms train exactly as E03's, it is the coloured-probe branch of
\cref{prop:sm:coloured} (\code{trace\_by\_arch}), which returns the same number at a
fifth of the cost. The checkpoint is selected on the implicit
objective $\E[\tfrac12|\widehat s|^2+\operatorname{div}\widehat s]$ of a held-out
noisy bundle, one view per validation sequence, with the exact trace
(\code{\_ism\_objective}, \code{select\_on="ism"}); the true score enters only in the
final scoring on the test bundle.
\end{definition}

In the paper's words (Section 5): (A) is ``clean sequences with a fresh corruption at
every step''; (B) is ``clean sequences with one fixed corruption each''; (C) is ``one
noised view per sequence and nothing else, which is what EM--BP sees'', and the
matched comparison adds that ``since the true score is not available for selection in
that setting, the checkpoint [is] selected on the implicit score-matching objective of
a held-out noisy bundle''. The sigma-algebras are nested: C sees strictly less than B,
which sees strictly less than A (A sees unboundedly many corruptions of each sequence
as training proceeds). The audit of 18--19 September 2026 (\code{docs/SCIENCE\_AUDIT.md})
is the reason the paper now names the interface of every number: an earlier draft
called the interface-B arm ``information-matched'', which it is not, because its
regression target is the clean sequence.

\why{The paper's headline factor at ``matched information'' is the ratio of the
interface-C network's error to the structured estimator's on the same bundles; it is a
statement about what the factorisation and the exact E-step are worth at equal
information, and \cref{def:sm:C} is what ``equal information'' means. Interface B
isolates the value of fresh corruptions, and the difference between A and B is that
value.}

\section{Worked example}
\label{sec:sm:example}

\begin{example}[The three objectives on a Gaussian]
\label{ex:sm:gaussian}
Let $\nsites=1$, $\Pz=\normal01$, so that $\Pt=\normal01$ at every $t$ and the score is
$\score(x)=-x$. Consider the linear family $s_c(x)=-cx$.
\begin{itemize}
\item Explicit: $J_{\mathrm{ESM}}(c)=\E(-cx+x)^2=(1-c)^2$.
\item Implicit: $J_{\mathrm{ISM}}(c)=\E[\tfrac12c^2x^2+(-c)]=\tfrac12c^2-c=\tfrac12(c-1)^2-\tfrac12$.
  The offset $-\tfrac12$ is $-\tfrac12\E|\nabla\log p|^2$, as \cref{eq:sm:ism} says,
  and the minimiser is $c=1$ in both.
\item Denoising: with $x=\alpha a+\sqrt{\Dt}z$ and target $-z/\sqrt{\Dt}$,
  $J_{\mathrm{DSM}}(c)=\E(-cx+z/\sqrt{\Dt})^2=c^2\E x^2-2c\,\E[xz]/\sqrt{\Dt}+\E z^2/\Dt
  =c^2-2c+1/\Dt=(c-1)^2+e^{-2t}/\Dt$, since $\E[xz]=\sqrt{\Dt}$ and $1/\Dt-1=e^{-2t}/\Dt$.
  The offset is the irreducible term $(e^{-2t}/\Dt^2)\mathrm{MMSE}_t$ of
  \cref{cor:sm:channel} with $\mathrm{MMSE}_t=\Var(a\mid x)=1-\alpha^2/(\alpha^2+\Dt)=\Dt$.
  At $t=\log2$ it equals $(1/4)/(3/4)=\tfrac13$, the SNR.
\end{itemize}
The \code{eps} loss for the same family is $\Dt J_{\mathrm{DSM}}$ and the \code{x0}
loss is $(\Dt^2/\alpha^2)J_{\mathrm{DSM}}$: at $t=\log2$ these are $\tfrac34J_{\mathrm{DSM}}$
and $\tfrac94J_{\mathrm{DSM}}$, and at level 1 of \cref{tab:ou:levels} they are
$0.095\,J_{\mathrm{DSM}}$ and $0.010\,J_{\mathrm{DSM}}$.
\end{example}

\where{Score matching is the training principle of every network in the paper's
Section 5 and Appendix C; the interfaces are defined in Section 5 and used in the
baselines and matched-comparison paragraphs; the estimand is Appendix C. Experiments:
E03 (screen, confirmatory table, matched arm, length extrapolation), E08
(misspecification network arm), E14 (regime baselines, whose network arms train
exactly as E03's: window head with the exact trace by coloured probes, transformer
with Hutchinson's estimator), and E05 for the network samplers. Code: \code{src/\allowbreak markov\_score/\allowbreak nn\_torch.py} (\code{TrainConfig},
\code{train}, \code{dsm\_\allowbreak loss}, \code{ism\_\allowbreak loss}, \code{divergence},
\code{ism\_\allowbreak objective}, \code{Denoiser}), \code{experiments/\allowbreak e03\_baselines.py}
(\code{\_\allowbreak train\_\allowbreak select}, \code{\_\allowbreak one\_\allowbreak view\_\allowbreak groups},
\code{\_\allowbreak train\_\allowbreak select\_\allowbreak fixed\_\allowbreak corruption}, \code{\_\allowbreak ism\_\allowbreak objective}, \code{\_\allowbreak risk},
\code{\_\allowbreak sq}), \code{experiments/\allowbreak protocol.py} (\code{estimand}, \code{bulk\_\allowbreak margin},
\code{net\_\allowbreak steps}), \code{tests/\allowbreak test\_nn\_torch.py}.}

\paragraph{References.}
A.~Hyv\"arinen, 2005, ``Estimation of non-normalized statistical models by score
matching'', Journal of Machine Learning Research 6, 695--709.
P.~Vincent, 2011, ``A connection between score matching and denoising autoencoders'',
Neural Computation 23, 1661--1674.
M.~F.~Hutchinson, 1990, ``A stochastic estimator of the trace of the influence matrix
for Laplacian smoothing splines'', Communications in Statistics: Simulation and
Computation 19, 433--450.
A.~R.~Curtis, M.~J.~D.~Powell and J.~K.~Reid, 1974, ``On the estimation of sparse
Jacobian matrices'', Journal of the Institute of Mathematics and its Applications 13,
117--119.
Y.~Song, S.~Garg, J.~Shi and S.~Ermon, 2019, ``Sliced score matching: a scalable
approach to density and score estimation'', Uncertainty in Artificial Intelligence.
Y.~Song, J.~Sohl-Dickstein, D.~P.~Kingma, A.~Kumar, S.~Ermon and B.~Poole, 2021,
``Score-based generative modeling through stochastic differential equations'',
International Conference on Learning Representations.
